\chapter{Non-Abelian Gauge Theory}

\section{Covariant derivative}

Let \(\psi\) transform as
\[
  \psi\to U(x)\psi,
  \qquad
  U(x)=e^{-\ii\alpha^a(x)T^a}.
\]
Define
\[
  D_\mu=\p_\mu+\ii g A_\mu,
  \qquad
  A_\mu=A_\mu^aT^a.
\]
Demand
\[
  D_\mu\psi\to U D_\mu\psi.
\]
This fixes the gauge-field transformation:
\[
  A_\mu'\ =
  U A_\mu U^{-1}
  +\frac{\ii}{g}(\p_\mu U)U^{-1}.
\]
For infinitesimal \(U=1-\ii\alpha\), where \(\alpha=\alpha^aT^a\),
\[
  A_\mu'=A_\mu+\frac1g\p_\mu\alpha+\ii[A_\mu,\alpha]+O(\alpha^2).
\]
In components,
\[
  \delta A_\mu^a=\frac1g\p_\mu\alpha^a-f^{abc}\alpha^bA_\mu^c.
\]

\section{Field strength from a commutator}

Compute
\[
  [D_\mu,D_\nu]
  =
  [\p_\mu+\ii gA_\mu,\p_\nu+\ii gA_\nu].
\]
Acting on a test field,
\[
  [D_\mu,D_\nu]
  =
  \ii g(\p_\mu A_\nu-\p_\nu A_\mu)
  -g^2[A_\mu,A_\nu].
\]
Write this as
\[
  [D_\mu,D_\nu]=\ii g F_{\mu\nu},
\]
where
\[
  F_{\mu\nu}
  =
  \p_\mu A_\nu-\p_\nu A_\mu+\ii g[A_\mu,A_\nu].
\]
Using \([T^a,T^b]=\ii f^{abc}T^c\),
\[
  F_{\mu\nu}^a
  =
  \p_\mu A_\nu^a-\p_\nu A_\mu^a
  -g f^{abc}A_\mu^bA_\nu^c.
\]
Sign conventions vary with the sign chosen in \(D_\mu\); physical predictions do not.

\section{Yang--Mills action}

The gauge-invariant Lagrangian is
\[
  \Lag_{\text{YM}}=-\frac12\Tr(F_{\mu\nu}F^{\mu\nu}).
\]
With the common normalization
\[
  \Tr(T^aT^b)=\frac12\delta^{ab},
\]
this becomes
\[
  \Lag_{\text{YM}}=-\frac14F_{\mu\nu}^aF^{a\mu\nu}.
\]
Because \(F_{\mu\nu}\) contains \(A_\mu A_\nu\), the Lagrangian contains cubic and quartic gauge-boson interactions.  This is the major difference from electromagnetism: non-Abelian gauge bosons carry the charge of their own gauge group.

\section{Equations of motion}

Vary the Yang--Mills action.  The variation of the field strength is
\[
  \delta F_{\mu\nu}=D_\mu\delta A_\nu-D_\nu\delta A_\mu.
\]
Then
\[
  \delta\Lag
  =
  -\Tr(F^{\mu\nu}\delta F_{\mu\nu})
  =
  -2\Tr(F^{\mu\nu}D_\mu\delta A_\nu),
\]
using antisymmetry.  Integrating by parts with the covariant derivative gives
\[
  \delta S=
  2\int\dd^4x\,\Tr((D_\mu F^{\mu\nu})\delta A_\nu).
\]
Thus the source-free Yang--Mills equation is
\[
  D_\mu F^{\mu\nu}=0.
\]
With matter, it becomes
\[
  D_\mu F^{\mu\nu}=j^\nu.
\]

\section{Gauge fixing and ghosts}

In the path integral, integrating over all \(A_\mu\) counts physically equivalent gauge copies infinitely many times.  The Faddeev--Popov idea multiplies by an identity:
\[
  1=\Delta_{\text{FP}}[A]
  \int\D\alpha\,\delta(G(A^\alpha)),
\]
where \(G(A)=0\) is a gauge condition.  For covariant gauges,
\[
  G^a(A)=\p_\mu A^{a\mu}.
\]
The determinant \(\Delta_{\text{FP}}\) can be represented by anticommuting scalar fields \(c^a,\bar c^a\), called ghosts:
\[
  \Lag_{\text{ghost}}=\bar c^a(-\p^\mu D_\mu^{ab})c^b.
\]
In non-Abelian theory \(D_\mu^{ab}\) contains \(A_\mu\), so ghosts interact with gauge fields.  They are not external particles; they are bookkeeping fields required by gauge redundancy and unitarity in covariant gauges.

\section{BRST symmetry in one page}

After gauge fixing, ordinary gauge symmetry is hidden, but a fermionic symmetry remains: BRST symmetry.  With anticommuting parameter suppressed, the transformations have the form
\[
  sA_\mu^a=D_\mu^{ab}c^b,
  \qquad
  sc^a=-\frac12g f^{abc}c^bc^c,
  \qquad
  s\bar c^a=B^a,
  \qquad
  sB^a=0.
\]
Here \(B^a\) is an auxiliary field.  The operation \(s\) is nilpotent:
\[
  s^2=0.
\]
Physical states are selected by a cohomology condition: states closed under \(s\), modulo states that are themselves \(s\)-variations.  The formalism is compact because it packages gauge invariance, ghost cancellations, and unitarity into one algebraic statement.

\section{Asymptotic freedom}

For a Yang--Mills theory with gauge group \(SU(N)\) and \(n_f\) Dirac fermions in the fundamental representation, the one-loop beta function has the form
\[
  \mu\frac{\dd g}{\dd\mu}
  =
  -\frac{g^3}{16\pi^2}
  \left(\frac{11}{3}N-\frac{2}{3}n_f\right)+O(g^5).
\]
If the bracket is positive, the coupling becomes weaker at high energies.  This is asymptotic freedom.  The minus sign is a quantum effect of gauge-boson self-interactions overpowering matter screening.  Quantum chromodynamics, the theory of quarks and gluons, has this property.

\section*{Exercises}
\addcontentsline{toc}{section}{Exercises}

\begin{exercise}
Show that the Abelian limit \(f^{abc}=0\) removes the \(A_\mu A_\nu\) part of \(F_{\mu\nu}^a\).
\begin{answer}
The commutator of gauge fields is proportional to \(f^{abc}\), so it vanishes when the group is Abelian.
\end{answer}
\end{exercise}

\begin{exercise}
Explain why non-Abelian gauge bosons can interact with one another.
\begin{answer}
The field strength contains a term quadratic in \(A_\mu\); squaring \(F_{\mu\nu}\) in the Lagrangian produces cubic and quartic powers of \(A_\mu\).
\end{answer}
\end{exercise}
